\begin{proof}
\begin{enumerate}[label=\arabic*.]
\item We give the error bounds and the complete covering argument behind
the interval computation.  This is a finite proof of the minimum,
not an inference from the preliminary grid search.  All coefficients,
accepted rectangles and local continuation bounds are retained in the
four \texttt{certified\_minimum\_*} programs and their JSON outputs.
The file \texttt{minimum\_certificate\_summary.json} records their
SHA-256 digests as well as sharp interval endpoints.

\item Put $a=|\lambda_*|$.  The normalized root coordinates are
\[
 z_\infty=(a/\lambda)^{1/5},\qquad
 z_0=(\lambda/a)^{1/4}.
\]
In the regular root frames of P1.16 their connections have the form
\[
 G(z)=\frac{\sum_{j=0}^{m-1}B_jz^j}{1+z^m},\qquad m=5,4.
\tag{1.18.2}
\]
The entries of $B_j$ are the exact algebraic numbers obtained by
substituting the positive real number $a^{\pm1/m}$ in the rational
connection.  The generator of the finite deck group acts on these
frames by their stated eigencharacters.  Since its flat matrix is
real and fixes $e$, it leaves $h_b$ invariant.  It therefore suffices
to cover the sector $0\leq\arg z\leq2\pi/m$ in each root disc.

\item Use the positive diagonal weights
\[
 (1,1/128,1/128,1/32,1/64,1/32)
\]
and the corresponding weighted maximum row norm.  Exact ball
arithmetic gives
\[
 \left\|\sum_j |B_j|\right\|<1.447\quad(m=5),\qquad
 \left\|\sum_j |B_j|\right\|<1.135\quad(m=4).
\]
Here the absolute values are entrywise, before taking the row norm.
Let $B$ denote the actual enclosed bound.  The scalar inequality
\[
 \frac1{1-t^m}\leq\frac1{m(1-t)}+\frac{m-1}{m},\qquad0\leq t<1,
\]
follows on multiplying by $m(1-t)(1+\cdots+t^{m-1})$;
the remaining assertion is
$\sum_{j=1}^{m-1}t^j\geq(m-1)t^m$.
Gronwall's inequality consequently bounds the fundamental solution
$U(0)=I$ on $|z|\leq R<1$ by
\[
 K_R=\exp\!\left[B\left(-\frac{\log(1-R)}m
                    +\frac{m-1}{m}R\right)\right].
\]
Cauchy's estimate bounds the tail after degree $n$ on $|z|\leq r<R$
by
\[
 K_R\frac{(r/R)^{n+1}}{1-r/R}.
\tag{1.18.3}
\]
The code uses $R=0.999$.  The coefficients are computed directly
from (1.18.2):
\[
 nU_n=-(n-m)U_{n-m}
       +\sum_{j=0}^{\min(m-1,n-1)}U_{n-1-j}B_j,
\]
with the first term omitted when $n<m$.
There is no derivative approximation in this recursion.
The regular limit frame is determined by certified continuation
from $\lambda=1$ to a positive root-chart point and inversion of
this convergent series.  At the infinity center it yields
\[
 h_b(\infty)\in
 [-1.094540250311402780116268263054945704
       \ \mathbin{\pm}\ 3\cdot10^{-37}].
\tag{1.18.4}
\]
The corresponding enclosure at zero has center
\[
-0.9317933763565235054307951057960111149908282005165194125
\]
and radius $7.91\cdot10^{-56}$.

\item We explain two numerical details which are needed to make these
bounds rigorous and efficient.  At a root center, let $H$ denote
the actual limiting period frame, and form the real matrix whose
five columns are
\[
 \operatorname{Im}H_{\mathrm{low},1},\quad
 \operatorname{Re}H_{\mathrm{low},1},\quad
 \operatorname{Im}H_{\mathrm{low},2},\quad
 \operatorname{Re}H_{\mathrm{low},2},\quad
 \operatorname{Im}H_{\mathrm{low},3}.
\]
Its certified inverse gives a constant real change of the lower five
flat coordinates.  Subtracting the corresponding real linear
combination of those coordinates from the top coordinate also
makes the first two top entries zero at the center.  This is an
explicit real invertible change fixing $e$.  It does not change
$h_b$: on the real intersection vector the extra top term is real,
so its imaginary part is zero.  This change greatly improves the
four real equations used to evaluate (1.15.3).

\item For a polynomial $p(z)=\sum c_jz^j$ with ball coefficients, evaluation
on a rectangle centered at $z_c$, of enclosing radius $\rho$, uses
its midpoint-coefficient polynomial at the point $z_c$.  Coefficient
uncertainty is bounded by
$\sum_j\operatorname{rad}(c_j)|z_c|^j$, and variation over the
rectangle by
$\rho\sum_j j|c_j|r^{j-1}$, where $r$ bounds the modulus on it.
The tail (1.18.3) is then added.  The point evaluation uses 800-digit
ball arithmetic; its rounding radius is retained.  These explicit
bounds avoid the excessive interval enlargement of long rectangular
Horner recursions.  No error radius is discarded.

\item On the square $|\operatorname{Re}z_\infty|,
|\operatorname{Im}z_\infty|\leq0.1$, differentiating the exact linear
system for the real intersection vector gives the following
certified bounds for the real Hessian of $h_b$:
\[
\begin{split}
 0.2898644&<h_{xx},h_{yy}<0.3946911,\\
 |h_{xy}|&<0.0524256,\qquad
 \det\operatorname{Hess}(h_b)>0.0757786.
\end{split}
\tag{1.18.5}
\]
The derivatives of the period frame are evaluated by
$Y'=YG$ and $Y''=Y(G^2+G')$; the derivatives of the intersection
vector follow by differentiating its invertible four-by-four real
system.  Finite rotation invariance forces $dh_b(0)=0$ in this
root chart.  Positive definiteness on the convex square therefore
proves a strict minimum at its center.

\item Outside that square, a recursive dyadic rectangle cover proves
$h_b>h_b(\infty)$ throughout $|z_\infty|\leq0.98$ in its sector.
It has 900 accepted rectangles, maximum subdivision depth eight.
The zero-root disc $|z_0|\leq0.98$ is certified by 463 rectangles,
also of maximum depth eight.  A rectangle is discarded only when
its entire intersection with the required sector and disc is empty,
or when it lies in the already proved local square.  Every remaining
rectangle is either certified by the strict interval inequality or
subdivided.  The complete records show that the subdivision terminates;
there is no unresolved rectangle.

\item It remains to cover a cusp neighborhood and the collar between these
two root discs.  Write
\[
 q=(\lambda-\lambda_*)/\lambda_*.
\]
In this coordinate the connection is $R/q+A(q)$, with $R^2=0$
and $A$ analytic on $|q|<1$.  A regular-singular factorization is
\[
 Y(q)=H\exp((\log q)R)P(q),\qquad P(0)=I.
\]
The recursion for its analytic factor is
\[
 (n-\operatorname{ad}_R)P_n
   =\sum_{j=0}^{n-1}P_{n-1-j}A_j,
 \qquad \operatorname{ad}_R(Q)=QR-RQ.
\]
Because $R^2=0$, the cube of this adjoint operator vanishes and its
inverse is the finite expression
$1/n+\operatorname{ad}_R/n^2+\operatorname{ad}_R^2/n^3$.
In an exact dyadic diagonal norm the code proves
$\|R\|<0.841$ and uses the sharper recorded bound
\[
 C=1+2\|R\|+4\|R\|^2<5.508.
\]
On the circle of radius $0.4$, the absolute-coefficient majorant of
$A$ is less than $4.895$.  The coefficient recursion is dominated
by the scalar equation $p'=Ca_{\rm abs}p$; hence
$\|P\|\leq\exp(C\cdot4.895\cdot0.4)$ there.  Cauchy's estimate
now gives the same explicit geometric tail bound as (1.18.3), with
this exponential in place of $K_R$ and $R=0.4$.

\item The constant $H$ is obtained by certified continuation to $q=0.01$
along a specified polygon in the upper half-plane.  The matrix
$2\pi iHRH^{-1}$ is integral, since it is the logarithm of the
corresponding marked local monodromy.  Its ball enclosure contains
exactly one integer matrix; that matrix and an integral Jordan
marking $J$ are exported.  They satisfy
\[
 J^{-1}(2\pi iHRH^{-1})J=\begin{pmatrix}0&I_3\\0&0\end{pmatrix}.
\]
At the fixed parameter $c_*=-h_b(\infty)$, the last three columns of
$J^{-1}HP(q)$ specifying $F_{c_*}$ have blocks $T(q),D(q)$.
The lower block is certified invertible on the stated cusp disc.
The period matrix is therefore
\[
 Z(q)=\frac{\log q}{2\pi i}I_3+C(q),\qquad C(q)=T(q)D(q)^{-1}.
\]
Put $B(q)=\operatorname{Im}C(q)$ and $y=-\log|q|/(2\pi)$.
Then
\[
 \det\operatorname{Im}Z
 =y^3+\operatorname{tr}(B)y^2
   +\Bigl(\sum_{i<j}(B_{ii}B_{jj}-B_{ij}B_{ji})\Bigr)y+\det B.
\tag{1.18.6}
\]
On the whole disc $|q|\leq0.01$, interval substitution in (1.18.6),
followed by $y=-\log(0.01)/(2\pi)+v$, gives four strictly positive
coefficients in $v$; their certified lower bounds, in increasing
order, are $0.8111$, $3.6097$, $3.5346$, and $1$.
The sharp endpoints are in \texttt{cusp\_cubic\_bounds.json}.
Since $v\geq0$, the determinant is positive
throughout the punctured disc.  Near the cusp $h_b+c_*>0$ by P1.8;
continuity and transversality then give this same strict inequality
throughout that disc.  A separate finite cover of
$0.01\leq|q|\leq0.25$, using the same analytic factor and (1.18.6),
contains 7033 certified rectangles, of maximum depth seven.  It
extends the inequality to the whole punctured disc $|q|\leq0.25$.

\item The only remaining region, expressed in the infinity-root chart, is
\[
 0.98\leq |z_\infty|\leq0.98^{-4/5},\qquad
 0\leq\arg z_\infty\leq2\pi/5,
 \qquad |q|\geq0.25,
\]
where $q=-z_\infty^{-5}-1$.
Each rectangle in this collar is reached from radius $0.95$ by
Taylor continuation of (1.18.2).  The local analytic radius is
$0.449$ times the distance to the nearest root of $1+z^5$.
The numerator is bounded in the same diagonal norm and the
denominator by the product of the five distances to those roots.
Gronwall and Cauchy give an explicit tail for every step, exactly
as in P1.16.  There are 427 steps and 117 accepted rectangles,
of maximum depth eight; each proves $h_b>h_b(\infty)$.

\item These regions cover every point of the completed base except the
cusp itself, where $h_b$ tends to positive infinity.  Every inequality
away from the infinity center is strict.  This proves both the
global minimum and its uniqueness, and P1.8 gives the exact parameter
domain.

\item Finally, \texttt{normalization\_shift.py} evaluates
$a_{\rm split}=96Y_{03}(\lambda_*/2)/(625\pi i)$ by the same
certified continuation, with the polygon stated above.
The real and imaginary radii are respectively less than
$4.31\cdot10^{-56}$ and $3.15\cdot10^{-56}$.
Subtracting its imaginary part from (1.18.4) gives the stated
point-normalized minimum and interval.  This is the explicit
splitting morphism already established in P1.15, now numerically
certified.

\end{enumerate}
\end{proof}
